% =====================================================================
%  Avance de Gestión de Proyectos — Escanea y Entra (TESCHA)
%  Plantilla APA 7.ª edición (clase apa7)
%  Compilar con:  latexmk -pdf avance-niza.tex   (usa biber automáticamente)
% =====================================================================
\documentclass[stu,11pt,floatsintext,biblatex]{apa7}

% Fuente sans serif a 11 pt (APA permite sans serif de 11 pt, como Arial)
\usepackage{tgheros}
\renewcommand{\familydefault}{\sfdefault}

\usepackage[spanish,es-tabla]{babel} % es-tabla: "Tabla" en vez de "Cuadro"
\usepackage{csquotes}                % requerido por biblatex-apa
\addbibresource{citas.bib}

\usepackage{tikz}
\usetikzlibrary{arrows.meta,positioning,shapes.geometric}

% La clase apa7 deja el texto en bandera (raggedright) tras \maketitle;
% aquí se justifica, como pide el profesor.
\usepackage{ragged2e}
\usepackage{array}
% Columna de tabla con texto en bandera (para celdas largas).
\newcolumntype{P}[1]{>{\RaggedRight\arraybackslash}p{#1}}

% Logos arriba de la portada (los pide la escuela, no APA).
\newcommand{\hrm}{\noindent
  \includegraphics[height=1.6cm]{img/tescha_logo.png}\hfill
  \includegraphics[height=1.6cm]{img/logo_sitemas.png}\par}

% Índice después de la portada (lo pide la escuela; APA puro no lo usa).
\setcounter{tocdepth}{2}
\makeatletter
\patchcmd{\maketitle}{\@ifundefined{def@donotrepeattitle}}
  {\tableofcontents\newpage\@ifundefined{def@donotrepeattitle}}{}{}
\patchcmd{\maketitle}{\section*{\protect\normalfont\textbf{\@title}}}
  {\section*{\protect\normalfont\textbf{\@title}}\addcontentsline{toc}{section}{\@title}}{}{}
\makeatother
\AtBeginDocument{\defbibheading{bibliography}[\refname]{%
  \clearpage\section*{\normalfont\textbf{#1}}\addcontentsline{toc}{section}{#1}}}
\newcommand{\apendiceindice}[2]{\addcontentsline{toc}{section}{\appendixname~#1: #2}}

\fancypagestyle{titlepage}{\fancyhf{}\renewcommand{\headrulewidth}{0pt}}

\renewcommand{\figurenote}[1]{\par\noindent\textit{\notelabel.} #1\par}
\renewcommand{\tablenote}[1]{\par\smallskip\figurenote{#1}}

\hypersetup{hidelinks}
\setlength{\headheight}{15.4pt}

\DeclareDelimFormat[parencite]{finalnamedelim}{\addspace\bibstring{and}\space}
\DeclareDelimFormat[bib,biblist]{finalnamedelim}{%
  \ifthenelse{\value{listcount}>\maxprtauth}{}{\addspace\bibstring{and}\space}}
\DefineBibliographyExtras{spanish}{\protected\def\bibrangedash{\textendash}}

\AtBeginDocument{%
  \renewcommand{\keywordname}{Palabras clave}%
  \renewcommand{\acksname}{Nota del autor}%
  \renewcommand{\notelabel}{Nota}%
}

% ----------------------------- Datos ---------------------------------
\title{Escanea y Entra: Prototipo de Acceso QR para el TESCHA}
\shorttitle{ACCESO QR TESCHA}

\authorsnames{Juan Esteban Anzures Campos, Daniel Arellano Palacios, Ian Alexis Flores Martínez, Miguel Ángel Jurado Delgadillo, Juan Fernando Reyes Arguello, Bryan Uriel Villegas Raymundo}
\authorsaffiliations{{Ingeniería en Sistemas Computacionales, Tecnológico de Estudios Superiores de Chalco}}

% Modo stu
\course{Gestión de Proyectos de Software (grupo 4751, turno vespertino)}
\professor{Niza Lucero Alpízar Martínez}
\duedate{9 de octubre de 2026}

\makeatletter
\@ifundefined{def@stu}{}{\let\@shorttitle\@undefined \let\@acks\@undefined}
\makeatother

\begin{document}
\maketitle
\justifying

Este documento reúne, en un solo texto, lo que el equipo ha escrito hasta ahora
sobre el proyecto \textit{Escanea y Entra: Prototipo de Acceso QR para el
TESCHA}, visto desde la gestión de proyectos. Unifica la ficha inicial del
proyecto, el plan de calidad, el tablero de metas y seguimiento, la
presentación del avance y el documento de investigación del equipo.

El proyecto busca probar, en un solo acceso del plantel, si un código QR en el
celular del estudiante vuelve más rápida y más confiable la entrada. Las
secciones siguientes explican por qué se hace, qué se va a construir, hasta
dónde llega, quién hace cada parte, cuánto tiempo y dinero se prevé, cómo se
cuidará la calidad y cómo se dará seguimiento. Las partes técnicas se explican
con palabras sencillas.

Algunos datos todavía no existen porque el proyecto apenas arranca. Cuando
una cifra depende de las mediciones del diagnóstico o del piloto, el documento
la deja como pendiente y no la estima.

\section{Iniciación del Proyecto}
La iniciación responde qué se quiere resolver, para quién y con qué límites.
La Tabla~\ref{tab:ficha} resume la ficha inicial y los datos de autorización.

\begin{table}[tbp]
	\caption{Ficha Inicial del Proyecto Escanea y Entra}
	\label{tab:ficha}
	\small\setlength{\tabcolsep}{3pt}
	\begin{tabular}{P{3cm}P{12cm}}
		\toprule
		Campo & Contenido \\
		\midrule
		Nombre & Escanea y Entra: Prototipo de Acceso QR para el TESCHA. \\
		Necesidad o problema & El acceso de los estudiantes se controla con la revisión visual de credenciales físicas. Es lento en los minutos previos al inicio de clases, no deja un registro digital de quién entra ni a qué hora y no permite comprobar que quien muestra la credencial sea su titular. \\
		Propósito & Comprobar si un QR en el celular agiliza la entrada y deja constancia de cada ingreso. \\
		Patrocinio & Dirección del TESCHA. Docente responsable: Niza Lucero Alpízar Martínez. \\
		Usuarios & Estudiantes del turno vespertino (muestran su código QR), personal de vigilancia del acceso piloto (opera el acceso) y Departamento de Control Escolar (consulta los registros). \\
		Producto & Un torniquete a escala con un lector de código QR, un servidor que valida el código y guarda cada intento, una página donde el estudiante ve su código y un panel web de consulta. \\
		Valor esperado & Ingreso más rápido que con la revisión manual, menos ingresos con credenciales ajenas o capturas compartidas y un registro confiable de entradas para detectar ausentismo y planear el personal por horario. \\
		Enfoque de trabajo & Híbrido: etapas, hitos y fechas fijas de forma predictiva, y desarrollo del lector, el servidor y el panel en iteraciones semanales. \\
		Periodo & 17 semanas, del 5 de octubre de 2026 al 29 de enero de 2027, en un solo acceso peatonal del plantel. \\
		Autorización & Plan del proyecto entregado el 18 de septiembre de 2026. Arranque el 5 de octubre. Primer hito: permiso de la dirección (16 de octubre de 2026). \\
		Criterio de éxito & Que durante el piloto el tiempo promedio de validación por estudiante y la fila máxima sean menores que en la línea base, sin aceptar códigos vencidos o alterados y sin perder ningún registro. \\
		\bottomrule
	\end{tabular}
	\tablenote{Elaboración propia con base en la Ficha Inicial del Proyecto (Entregable 1) y la presentación del avance. TESCHA es el Tecnológico de Estudios Superiores de Chalco.}
\end{table}

\subsection{El Problema}
Hoy los vigilantes revisan las credenciales a simple vista. Eso tiene tres
consecuencias:
\begin{itemize}
	\item \textbf{Filas.} Se forman en los 15 a 20 minutos anteriores a la
	      primera clase del turno vespertino, cuando casi todos llegan en poco tiempo
	      y la revisión depende de una o dos personas.
	\item \textbf{Suplantación.} Una credencial prestada, vencida o ajena puede
	      pasar la revisión visual, porque nada comprueba que quien la muestra sea su
	      titular.
	\item \textbf{Falta de datos.} La institución no sabe cuántas personas entran
	      ni a qué hora, así que no puede detectar a tiempo a quienes dejan de asistir
	      ni planear cuántos accesos abrir y con cuánto personal.
\end{itemize}
Este problema es distinto del pase de lista en el salón. Un registro en la
entrada solo prueba que la persona entró al plantel, no que asistió a una
clase. Por eso el control de asistencia por materia queda fuera del proyecto.
Insatisfacciones parecidas con los registros manuales ya se han documentado en
otros planteles \parencite{torres2019qr}.

\subsection{Interesados}
Los usuarios usan el producto de forma directa. Los demás interesados se ven
afectados por el proyecto o pueden influir en él (Tabla~\ref{tab:interesados}).

\begin{table}[tbp]
	\caption{Usuarios e Interesados}
	\label{tab:interesados}
	\small\setlength{\tabcolsep}{3pt}
	\begin{tabular}{P{3.2cm}P{2.3cm}P{4.4cm}P{3.6cm}P{1.5cm}}
		\toprule
		Interesado & Tipo & Cómo participa & Qué espera & Influencia \\
		\midrule
		Estudiantes del turno vespertino & Usuario directo & Muestran su código QR desde el celular y responden la encuesta del diagnóstico. & Entrar rápido, sin filas, y que sus datos estén protegidos. & Media \\
		Personal de vigilancia del acceso piloto & Usuario directo & Opera el acceso junto al prototipo, atiende a quien no puede mostrar su código y consulta el panel. & Menos carga en la hora pico y una herramienta que detecte credenciales ajenas. & Media \\
		Departamento de Control Escolar & Usuario del panel & Proporciona la matrícula y consulta los reportes de ingreso. & Datos para detectar ausentismo y riesgo de deserción. & Alta \\
		Dirección del TESCHA & Patrocinador & Autoriza el piloto y decide si el sistema se extiende a otros accesos. & Evidencia de que la solución funciona y es de bajo costo. & Alta \\
		Área responsable de datos personales del TESCHA & Interesado & Revisa el aviso de privacidad y las condiciones para tratar datos de estudiantes. & Que el sistema cumpla la legislación de protección de datos. & Alta \\
		Docentes asesores & Interesado & Revisan los avances y orientan al equipo. & Entregables a tiempo y bien fundamentados. & Alta \\
		Equipo del proyecto & Ejecutor & Diseña, construye, prueba y documenta el prototipo. & Terminar el prototipo y aprobar la asignatura. & Alta \\
		\bottomrule
	\end{tabular}
	\tablenote{Fuente: Ficha Inicial del Proyecto (Entregable 1).}
\end{table}

\subsection{Supuestos y Restricciones}
El plan se apoya en dos supuestos: en el acceso hay toma de corriente y Wi-Fi,
y la mayoría de los estudiantes trae celular. Si alguno falla, el proyecto lo
trata como riesgo (sección~\ref{sec:riesgos}).

Las restricciones son un solo acceso peatonal, el turno vespertino, 17 semanas
con fechas que no se mueven y el presupuesto del propio equipo. El calendario
escolar no se mueve y el piloto debe terminar antes de las vacaciones de
invierno, mientras los estudiantes todavía asisten.

\subsection{Autorización y Criterio de Éxito}
El proyecto arranca con el plan entregado el 18 de septiembre de 2026 y con el
permiso de la dirección como primer hito, previsto para el 16 de octubre. Sin
ese permiso no se puede entrar al acceso. El criterio de éxito general es
completar el piloto y comparar tres indicadores contra la línea base. Los
criterios de éxito concretos están en la Tabla~\ref{tab:exito}.

\begin{table}[tbp]
	\caption{Criterios de Éxito}
	\label{tab:exito}
	\small\setlength{\tabcolsep}{3pt}
	\begin{tabular}{P{3.3cm}P{4.6cm}P{3.9cm}P{2.6cm}}
		\toprule
		Criterio & Meta & Cómo se verifica & Cuándo \\
		\midrule
		Tiempo promedio de validación por estudiante & Menor durante el piloto que en la línea base. & Hoja de observación del flujo en el acceso. & Diagnóstico y piloto \\
		Longitud máxima de la fila en el horario pico & Menor durante el piloto que en la línea base. & Hoja de observación del flujo en el acceso. & Diagnóstico y piloto \\
		Rechazo de códigos inválidos & Se rechaza el 100\,\% de los códigos vencidos, reenviados o alterados. & Casos de prueba en el salón. & Pruebas en el salón \\
		Integridad del registro & El 100\,\% de los intentos queda registrado; la diferencia con el conteo manual es de 2\,\% o menos. & Base de datos frente a la hoja de observación. & Piloto \\
		Aceptación & Al menos 85\,\% de los estudiantes del piloto pasa al primer intento sin ayuda. & Observación y registros del sistema. & Piloto \\
		Cumplimiento del cronograma & Piloto terminado el 18 de diciembre y entrega final el 29 de enero. & Seguimiento semanal del cronograma. & Todo el proyecto \\
		Decisión sustentada & Informe final con una recomendación respaldada por datos sobre extender el sistema a otros accesos. & Revisión del informe por los docentes asesores. & Cierre \\
		\bottomrule
	\end{tabular}
	\tablenote{Fuente: Ficha Inicial del Proyecto (Entregable 1). La línea base es la medición del acceso actual, que se levanta en el diagnóstico de octubre porque hoy no existe ninguna.}
\end{table}

\section{Objetivos del Proyecto}
\subsection{Objetivo General}
\textbf{Evaluar el efecto de un prototipo de control de acceso con códigos QR
	dinámicos sobre el tiempo de validación, la fila y las incidencias de
	identidad en el acceso peatonal piloto del TESCHA, para aportar evidencia que
	permita decidir si conviene ampliarlo a otros accesos.}

Hoy no existen datos que midan el problema, y una decisión sobre el sistema
debe apoyarse en evidencia. El objetivo pide evaluar el efecto, no demostrar
una mejora: si las mediciones no muestran diferencia, ese también es un
resultado válido.

\subsection{Indicadores}
El efecto se mide con tres indicadores, primero antes del piloto (línea base)
y después durante él, en el horario pico del turno vespertino:
\begin{itemize}
	\item \textbf{Tiempo de validación:} promedio por estudiante, desde que
	      llega al punto de revisión hasta que lo cruza.
	\item \textbf{Fila máxima:} la mayor cantidad de personas esperando en el
	      horario pico.
	\item \textbf{Incidencias de identidad:} credenciales ajenas detectadas.
\end{itemize}
La pregunta que guía el trabajo es: ¿cuál es el efecto de implementar un
sistema automatizado de validación mediante códigos QR dinámicos en la
eficiencia general del control de acceso peatonal de los estudiantes del
TESCHA? Por eficiencia se entiende dejar pasar con rapidez a quienes tienen
derecho a entrar, y solo a ellos.

\subsection{Objetivos Específicos}
Los ocho objetivos específicos se reparten en dos líneas de trabajo
(Tabla~\ref{tab:especificos}).

\begin{table}[tbp]
	\caption{Objetivos Específicos}
	\label{tab:especificos}
	\small\setlength{\tabcolsep}{3pt}
	\begin{tabular}{P{1.2cm}P{7.3cm}P{6.6cm}}
		\toprule
		Clave & Qué se hará & Por qué \\
		\midrule
		\multicolumn{3}{l}{\textit{Desarrollo tecnológico}} \\
		OT1 & Diseñar la arquitectura del sistema: lector, torniquete, servidor, base de datos y panel web. & El acceso actual no deja registro digital y el prototipo necesita una estructura que lo registre. \\
		OT2 & Implementar la generación y validación de códigos QR que cambian cada pocos segundos y están firmados, con pruebas de códigos vencidos, repetidos y alterados. & Un código fijo puede copiarse y compartirse; la firma resiste el reenvío y la modificación del mensaje \parencite{subpratatsavee2015hmac}. \\
		OT3 & Construir el prototipo funcional que integre lector, torniquete, servidor y panel web. & Sin prototipo no hay sistema que evaluar. \\
		OT4 & Verificar el desempeño técnico: tiempo de lectura y de consulta, y lecturas fallidas con distintas luces, brillos de pantalla y distancias. & No se encontró literatura que reporte el tiempo de lectura de un QR en este lector, así que debe medirse con el prototipo. \\
		\multicolumn{3}{l}{\textit{Investigación y análisis de datos}} \\
		OI1 & Analizar la situación actual del acceso con observación de cinco días, encuesta y entrevistas, para fijar la línea base. & No existe ningún registro del flujo de acceso; sin línea base no se puede atribuir ningún cambio al sistema. \\
		OI2 & Estimar la viabilidad y la aceptación entre los estudiantes: celular propio, conexión, batería y disposición a mostrar un QR. & El sistema depende del celular del estudiante y se desconoce si la comunidad lo aceptaría. \\
		OI3 & Interpretar las obligaciones de protección de datos personales aplicables al TESCHA \parencite{edomex2017datos, lgpdppso2025}. & El TESCHA es sujeto obligado de ambas leyes y el prototipo registra la hora de ingreso de cada estudiante. \\
		OI4 & Evaluar el efecto del sistema en los tres indicadores, comparando la línea base con el piloto. & La pregunta pide el efecto, y eso solo se conoce comparando un antes y un después. \\
		\bottomrule
	\end{tabular}
	\tablenote{Elaboración propia con base en el documento de investigación del equipo. OT es desarrollo tecnológico; OI es investigación y análisis de datos.}
\end{table}

Los objetivos de desarrollo producen el sistema que se evalúa. Los de
investigación miden los indicadores y aportan lo necesario para interpretarlos.
Sin los primeros no hay sistema que evaluar; sin los segundos no se sabe si el
sistema sirvió.

\section{Descripción de la Solución}
La solución es un torniquete a escala que se abre solo cuando el estudiante
muestra un código válido desde su celular. El código es un cuadro de puntos
llamado código QR, como los que se leen con la cámara del teléfono.

\subsection{Cómo Funciona, Paso a Paso}
\begin{enumerate}
	\item El estudiante abre una página web desde su celular y ve su código QR.
	      El código cambia cada pocos segundos.
	\item Al llegar al torniquete, muestra la pantalla a una cámara pequeña
	      (el lector ESP32-CAM), que captura la imagen y lee el código.
	\item Una segunda placa electrónica (el ESP32) recibe el código y lo envía
	      por Wi-Fi al servidor, que es la computadora central del sistema.
	\item El servidor revisa si el código es auténtico y sigue vigente, y
	      responde si el estudiante puede pasar.
	\item Si la respuesta es afirmativa, el torniquete se abre y una luz verde
	      lo indica. Si no, una luz y un zumbador avisan que no puede pasar.
	\item El servidor guarda cada intento, aceptado o rechazado, con fecha y
	      hora en una base de datos, que es el archivo donde se almacena la
	      información.
	\item El personal de vigilancia y Control Escolar consultan esos registros
	      en un panel web: ingresos por hora, intentos rechazados y estudiantes
	      que llevan varios días sin entrar.
\end{enumerate}

\begin{figure}[tbp]
	\caption{Cómo Funciona el Torniquete}
	\label{fig:flujo}
	\centering
	\begin{tikzpicture}[
		node distance=0.9cm and 1.0cm,
		box/.style={draw, rounded corners, align=center, minimum height=1.2cm,
				minimum width=2.4cm, font=\footnotesize},
		arr/.style={-{Latex[length=2mm]}, thick}
		]
		\node[box] (cel) {Celular\\(muestra el QR)};
		\node[box, right=of cel] (cam) {ESP32-CAM\\(lee el código)};
		\node[box, right=of cam] (esp) {ESP32\\(consulta y abre)};
		\node[box, below=of esp] (srv) {Servidor y\\base de datos\\(validan y guardan)};
		\node[box, left=of srv] (panel) {Panel web\\(vigilancia y\\Control Escolar)};
		\draw[arr] (cel) -- (cam);
		\draw[arr] (cam) -- (esp);
		\draw[arr] (esp) -- (srv) node[midway, right, font=\scriptsize]{Wi-Fi};
		\draw[arr] (srv) -- (panel) node[midway, above, font=\scriptsize]{consulta};
	\end{tikzpicture}
	\figurenote{Elaboración propia. El servidor responde si el código es válido y el ESP32 libera el torniquete; cada intento queda registrado en la base de datos.}
\end{figure}

\subsection{Por Qué el Código Cambia}
Cada código lleva una firma digital llamada HMAC, una contraseña de un solo
uso que solo el servidor sabe calcular. Si alguien cambia el código, la firma
deja de coincidir y se rechaza. Como además el código se renueva cada pocos
segundos, una captura de pantalla compartida deja de servir al poco tiempo.
Este tipo de firma resiste el reenvío y la modificación del mensaje
\parencite{subpratatsavee2015hmac}, y la renovación periódica ya se probó en
una universidad del Reino Unido, donde un código renovado cada 20 segundos
impidió que los estudiantes se reenviaran capturas \parencite{nwabuwe2023dynamicqr}.

\subsection{Límite de la Solución}
El código que cambia impide reutilizar capturas, pero no impide que un
estudiante preste su celular a otra persona. Por eso el proyecto busca reducir
la suplantación, no eliminarla. Además, como la medición se hace en un solo
acceso y sin grupo de comparación, los resultados se tomarán como evidencia
orientativa, no como prueba definitiva.

\section{Gestión del Alcance}
El alcance dice qué se hará y qué no. Se escribió desde el inicio para que, si
algo se atrasa, el equipo sepa qué puede ceder (Tabla~\ref{tab:alcance}).

\begin{table}[tbp]
	\caption{Alcance del Proyecto}
	\label{tab:alcance}
	\small\setlength{\tabcolsep}{3pt}
	\begin{tabular}{P{7.4cm}P{7.4cm}}
		\toprule
		Incluye & Queda fuera \\
		\midrule
		Un acceso peatonal piloto, en el turno vespertino. & Reconocimiento facial y tarjetas NFC. \\
		QR que cambia en el celular del estudiante, con su generación y validación. & Pase de lista dentro de los salones (asistencia por materia). \\
		Lector con cámara y torniquete a escala. & Registro de salidas y control de visitantes. \\
		Servidor que valida y registra cada intento en una base de datos. & Instalación en todos los accesos y torniquete industrial. \\
		Panel web de consulta para el personal. & Retirar la revisión manual durante el piloto. \\
		Medición del acceso antes y durante el piloto. & \\
		\bottomrule
	\end{tabular}
	\tablenote{Fuente: Ficha Inicial del Proyecto y presentación del avance.}
\end{table}

Durante el piloto la revisión manual sigue funcionando. Quien no pueda
mostrar su código entra como hasta ahora.

\subsection{Entregables}
Al cierre, el equipo entregará los productos de la Tabla~\ref{tab:entregables}.
Ninguno se da por terminado sin pasar su criterio de aceptación, y cada uno
tiene un responsable de la tabla de roles (Tabla~\ref{tab:roles}).

\begin{table}[tbp]
	\caption{Entregables del Proyecto}
	\label{tab:entregables}
	\small\setlength{\tabcolsep}{3pt}
	\begin{tabular}{P{1cm}P{7cm}}
		\toprule
		N.º & Entregable \\
		\midrule
		1 & Informe de diagnóstico \\
		2 & Prototipo del torniquete con lector \\
		3 & Servidor y base de datos \\
		4 & Página del estudiante con su QR \\
		5 & Panel web de consulta \\
		6 & Código fuente en un repositorio \\
		7 & Manual de uso y aviso de privacidad \\
		8 & Informe de resultados: antes frente a durante el piloto \\
		\bottomrule
	\end{tabular}
	\tablenote{Fuente: presentación del avance. La ficha inicial menciona los mismos productos: prototipo, informe de diagnóstico, código fuente, manual de uso, propuesta de aviso de privacidad e informe final.}
\end{table}

\section{Planificación del Proyecto}
El plan del proyecto se compone de seis planes: alcance, cronograma, costos,
calidad, riesgos, y recursos y comunicación. El alcance ya se describió; la
calidad tiene su propia sección (sección~\ref{sec:calidad}). Aquí se
presentan los demás.

\subsection{Triple Restricción}
Alcance, tiempo y costo dependen entre sí: si una se mueve, las otras dos lo
resienten.
\begin{itemize}
	\item \textbf{Alcance:} un acceso peatonal, turno vespertino y prototipo a
	      escala. Lo que queda fuera está escrito desde el inicio.
	\item \textbf{Tiempo:} 17 semanas fijas, del 5 de octubre de 2026 al 29 de
	      enero de 2027. El calendario escolar no se mueve.
	\item \textbf{Costo:} cerca de 4,100 pesos aportados por el propio equipo,
	      sin financiamiento externo, con una reserva de 10\,\%.
\end{itemize}
La regla del equipo es que, si algo se atrasa, se cede alcance antes que fecha
o calidad.

\subsection{Estrategias Generales}
Seis decisiones guían el proyecto:
\begin{itemize}
	\item \textbf{Empezar pequeño:} un acceso, un turno y un prototipo a escala.
	      Primero evidencia barata; escalar solo si los números la respaldan.
	\item \textbf{Trabajar en paralelo:} el hardware y el software avanzan
	      mientras se levanta el diagnóstico; la integración se deja para el final.
	\item \textbf{Usar lo barato y lo abierto:} componentes de bajo costo y
	      herramientas libres, en lugar de biometría o torniquetes industriales.
	\item \textbf{Medir antes y después:} el mismo instrumento se aplica antes
	      y durante el piloto, para comparar con números y no con impresiones.
	\item \textbf{Guardar holgura:} reserva del 10\,\% en costo y una semana
	      libre antes de vacaciones para absorber atrasos.
	\item \textbf{No retirar nada:} la revisión manual sigue funcionando todo el
	      piloto. El sistema se suma, no sustituye a nadie.
\end{itemize}

\subsection{Enfoque y Cronograma}
El enfoque es híbrido. Es predictivo en la estructura porque hay fechas que no
se pueden mover: el permiso de la dirección es requisito para entrar al
acceso, los materiales tardan en llegar, el piloto debe terminar antes de las
vacaciones y la entrega final es el 29 de enero de 2027. Es ágil en el
desarrollo porque aún no se sabe cuánto tarda el lector en leer un código ni
cómo afectan la luz, el brillo de la pantalla o la distancia. Esa
incertidumbre exige probar, medir y ajustar. El lector, el servidor y el panel
se construyen en iteraciones de una semana: cada lunes se eligen tareas de una
lista priorizada y cada viernes se muestra lo terminado y se revisa el avance
contra el cronograma. La Tabla~\ref{tab:etapas} muestra las etapas.

\begin{table}[tbp]
	\caption{Etapas, Forma de Trabajo e Hitos}
	\label{tab:etapas}
	\small\setlength{\tabcolsep}{3pt}
	\begin{tabular}{P{2.2cm}P{3.3cm}P{6.4cm}P{3cm}}
		\toprule
		Etapa & Fechas & Forma de trabajo & Hito \\
		\midrule
		Planeación & 5 al 16 de octubre & Predictiva: permisos, aviso de privacidad y lista de materiales. & Permiso de la dirección \\
		Diagnóstico & 12 al 30 de octubre & Predictiva: encuesta, entrevistas y observación del acceso durante cinco días; compra de materiales. & Informe de diagnóstico \\
		Desarrollo & 26 de octubre al 27 de noviembre & Ágil: cinco iteraciones semanales en dos frentes paralelos, hardware y lector, y servidor, base de datos y página web. & Prototipo integrado \\
		Pruebas & 23 de noviembre al 18 de diciembre & Iterativa en el salón (23 de noviembre al 4 de diciembre) y predictiva en el piloto (7 al 18 de diciembre). & Fin del piloto \\
		Cierre & 4 al 29 de enero & Predictiva: comparación de resultados, informe final y manual de uso. & Entrega final \\
		\bottomrule
	\end{tabular}
	\tablenote{Fuente: Ficha Inicial del Proyecto (Entregable 1). Las semanas del 21 y del 28 de diciembre son vacaciones de invierno y no tienen trabajo programado. Las tareas se registran en un tablero compartido y el código en un repositorio Git.}
\end{table}

Las actividades corren en cuatro frentes en paralelo y una integración al
final (Tabla~\ref{tab:frentes}). La integración se hace en el salón y, hasta
que pasa ahí, el prototipo no se instala en el acceso.

\begin{table}[tbp]
	\caption{Frentes de Trabajo}
	\label{tab:frentes}
	\small\setlength{\tabcolsep}{3pt}
	\begin{tabular}{P{5cm}P{2.6cm}P{7cm}}
		\toprule
		Frente y responsable & Semanas & Resultado \\
		\midrule
		Hardware y torniquete (Daniel) & 3 a 6 & Prototipo a escala con el lector montado. \\
		Lector y firma del código (Ian y Miguel) & 3 a 7 & Lector con cámara y código que caduca en segundos. \\
		Servidor, datos y panel (Juan Fernando) & 3 a 8 & Validación, registro y consultas del personal. \\
		Campo y pruebas (Bryan) & 2 a 5 y 10 a 11 & Diagnóstico, pruebas y piloto. \\
		Integración en el salón & 8 y 9 & Prototipo completo probado fuera del acceso. \\
		\bottomrule
	\end{tabular}
	\tablenote{Fuente: presentación del avance. La semana 1 es la del 5 de octubre de 2026.}
\end{table}

\subsection{Costos}
El presupuesto es una estimación preliminar del equipo, hecha con precios de
referencia del mercado mexicano. Se confirma al cotizar (Tabla~\ref{tab:costos}).

\begin{table}[tbp]
	\caption{Presupuesto Estimado}
	\label{tab:costos}
	\small\setlength{\tabcolsep}{3pt}
	\begin{tabular}{P{8cm}r}
		\toprule
		Categoría & Estimado (MXN) \\
		\midrule
		Hardware del prototipo & 1,750 \\
		Servidor, dominio y hosting (4 meses) & 970 \\
		Impresiones y material de campo & 600 \\
		Traslados y compras & 400 \\
		Reserva de contingencia (10\,\%) & 380 \\
		\midrule
		Total estimado & 4,100 \\
		\bottomrule
	\end{tabular}
	\tablenote{Fuente: presentación del avance. Las cerca de 800 horas de trabajo del equipo no se monetizan. Si la escuela presta el servidor, el total baja 970 pesos.}
\end{table}

El control de costos sigue cuatro reglas: una línea base aprobada antes de la
primera compra, cada compra registrada con comprobante y responsable, un umbral
de variación de $\pm$10\,\% (pasarlo obliga a replantear) y la comparación del
gasto acumulado contra lo planeado en la junta semanal.

\subsection{Recursos}
Con lo siguiente cuenta el equipo para ejecutar el proyecto:
\begin{itemize}
	\item \textbf{Humanos:} seis integrantes, unas 8 horas por semana cada uno,
	      cerca de 800 horas en total.
	\item \textbf{Materiales:} lector con cámara, placa de control, motor, luces
	      y zumbador, fuente de energía, base de MDF y cableado.
	\item \textbf{Técnicos:} servidor con base de datos, repositorio de código y
	      entorno de programación.
	\item \textbf{Institucionales:} permiso de la dirección, espacio con
	      corriente y Wi-Fi, y apoyo de vigilancia.
	\item \textbf{Financieros:} aportación del equipo, cerca de 4,100 pesos.
	\item \textbf{De información:} padrón de estudiantes de Control Escolar y
	      resultados del diagnóstico.
\end{itemize}
Lo que depende de terceros (permiso, Wi-Fi, padrón) es lo que encabeza el
registro de riesgos.

\subsection{Comunicación}
El equipo se comunica con tres mecanismos: una junta semanal de 20 minutos
para revisar el avance contra el cronograma, un reporte quincenal a la
docente con lo hecho y lo pendiente, y un tablero de tareas en el repositorio
con una columna por frente. Todo cambio que toque fecha, costo o alcance lo
autoriza el líder del proyecto y queda por escrito.

\subsection{Riesgos}
\label{sec:riesgos}
El registro tiene siete riesgos. Cada uno tiene un responsable que lo vigila
y se revisa en la junta semanal. Los dos primeros condicionan la fecha del
piloto (Tabla~\ref{tab:riesgos}).

\begin{table}[tbp]
	\caption{Registro de Riesgos}
	\label{tab:riesgos}
	\small\setlength{\tabcolsep}{3pt}
	\begin{tabular}{P{4.8cm}P{2.3cm}P{1.8cm}P{5.7cm}}
		\toprule
		Riesgo & Probabilidad & Impacto & Respuesta \\
		\midrule
		La dirección no autoriza el piloto & Media & Alto & Solicitud formal en la semana 1; plan B en el laboratorio. \\
		Wi-Fi inestable en el acceso & Alta & Alto & Modo fuera de línea con lista local y sincronización. \\
		Retraso en la entrega de componentes & Media & Medio & Compra con dos semanas de anticipación y repuestos. \\
		Mala lectura con sol directo & Media & Medio & Visera, ajuste de exposición y pruebas con varias luces. \\
		Estudiantes sin celular o sin batería & Alta & Bajo & El carril manual sigue abierto durante todo el piloto. \\
		Manejo indebido de datos personales & Baja & Alto & Aviso de privacidad, datos mínimos, cifrado y borrado. \\
		Vacaciones y carga de fin de semestre & Alta & Medio & Piloto cerrado antes de las vacaciones; enero para el informe. \\
		\bottomrule
	\end{tabular}
	\tablenote{Fuente: presentación del avance. Los nombres de los responsables de cada riesgo no aparecen en las fuentes.}
\end{table}

\subsection{Cierre del Proyecto}
El 29 de enero de 2027 se hará un cierre en seis pasos: verificar cada
entregable contra su criterio de aceptación, entregar el informe final con la
comparación de los tres indicadores, firmar un acta de aceptación por la
dirección y la docente, transferir el manual y dar una capacitación breve a
vigilancia y Control Escolar, borrar o resguardar los registros según el aviso
de privacidad, y documentar las lecciones aprendidas con una recomendación
sobre llevar el sistema a los demás accesos o detenerlo.

\section{Organización del Equipo}
Cada integrante tiene un rol principal (Tabla~\ref{tab:roles}). Eso no
significa que trabajen por separado: en el diagnóstico y en el piloto los seis
aplican encuestas y observan la entrada, porque la hoja de observación pide
dos observadores por día.

\begin{table}[tbp]
	\caption{Roles y Responsabilidades del Equipo}
	\label{tab:roles}
	\small\setlength{\tabcolsep}{3pt}
	\begin{tabular}{P{3.4cm}P{3.2cm}P{8.3cm}}
		\toprule
		Integrante & Rol & Responsabilidades \\
		\midrule
		Juan Esteban Anzures Campos (JEAC) & Líder del proyecto & Coordinar al equipo y dar seguimiento al cronograma, tratar con la dirección y los profesores, gestionar los permisos, reunir la documentación y los informes, y autorizar los cambios. \\
		Daniel Arellano Palacios (DAP) & Responsable de hardware & Elegir y comprar los materiales, armar el torniquete, conectar el motor, las luces y la alimentación, e instalar el prototipo en el acceso. \\
		Ian Alexis Flores Martínez (IAFM) & Responsable del lector & Programar el lector y la placa de control para leer el código, consultarlo con el servidor y abrir el torniquete; medir cuánto tarda cada lectura. \\
		Miguel Ángel Jurado Delgadillo (MAJD) & Responsable del servidor y la seguridad & Construir el servidor que genera y valida los códigos, proteger la comunicación entre las partes y preparar el aviso de privacidad junto con el líder. \\
		Juan Fernando Reyes Arguello (JFRA) & Responsable de datos y página web & Diseñar la base de datos de ingresos, la página donde el estudiante ve su código y el panel de consulta con los reportes. \\
		Bryan Uriel Villegas Raymundo (BUVR) & Responsable de campo y pruebas & Organizar la encuesta, las entrevistas y la observación de la entrada, capturar y analizar los datos, y coordinar las pruebas antes y durante el piloto. \\
		\bottomrule
	\end{tabular}
	\tablenote{Fuente: Ficha Inicial del Proyecto (Entregable 1). Cada función responde por sus entregables y por los riesgos de su frente.}
\end{table}

La matriz de la Tabla~\ref{tab:raci} indica, para cada actividad, quién la
realiza (R), quién rinde cuentas y aprueba (A), a quién se consulta (C) y a
quién se informa (I). Cada actividad tiene un solo responsable de rendir
cuentas.

\begin{table}[tbp]
	\caption{Matriz de Responsabilidades (RACI)}
	\label{tab:raci}
	\small\setlength{\tabcolsep}{3pt}
	\begin{tabular}{P{6.6cm}cccccc}
		\toprule
		Actividad & JEAC & DAP & IAFM & MAJD & JFRA & BUVR \\
		\midrule
		Planeación, cronograma y seguimiento semanal & A/R & C & C & C & C & C \\
		Permiso de la dirección & A/R & I & I & I & I & C \\
		Aviso de privacidad & A & I & I & R & C & I \\
		Encuesta, entrevistas y observación (diagnóstico) & C & R & R & R & R & A/R \\
		Informe de diagnóstico & A & I & I & I & C & R \\
		Compra de materiales & C & A/R & C & I & I & I \\
		Armado del torniquete & I & A/R & C & I & I & I \\
		Programación del lector & I & C & A/R & C & I & I \\
		Servidor y códigos QR dinámicos & I & I & C & A/R & C & I \\
		Base de datos, página del estudiante y panel web & I & I & I & C & A/R & C \\
		Integración y pruebas en el salón & I & R & R & R & R & A/R \\
		Instalación y piloto en el acceso & A & R & R & C & C & R \\
		Comparación de resultados & C & I & I & I & C & A/R \\
		Informe final y manual de uso & A/R & C & C & C & R & R \\
		\bottomrule
	\end{tabular}
	\tablenote{Fuente: Ficha Inicial del Proyecto (Entregable 1). R = realiza; A = rinde cuentas y aprueba; C = consultado; I = informado.}
\end{table}

\section{Plan de Calidad}
\label{sec:calidad}
El plan define qué significa calidad para el prototipo y cómo se comprobará.
Fija cinco objetivos de calidad y, para cada uno, criterios de aceptación
medibles, pruebas, responsable y métricas. Las características de calidad se
toman del modelo de calidad de producto de la norma ISO/IEC 25010. El plan
cubre las cinco partes del producto (página del estudiante, lector, torniquete,
servidor con base de datos y panel web) y la medición del acceso antes y durante
el piloto, de la que depende saber si el prototipo cumplió su propósito.

En la gestión de la calidad el equipo hace tres cosas. Planifica: escribe los
requisitos funcionales, de calidad y legales antes de construir, cada uno con
su criterio de aceptación y su responsable. Asegura: revisa el código entre
compañeros, usa una lista de verificación en el armado, hace un ensayo completo
en el salón y revisa el aviso de privacidad. Controla: hace pruebas de lectura,
del panel y del registro, y verifica cada entregable contra su criterio antes
de darlo por terminado.

\subsection{Objetivos de Calidad}
\begin{table}[tbp]
	\caption{Objetivos de Calidad del Prototipo}
	\label{tab:oc}
	\small\setlength{\tabcolsep}{3pt}
	\begin{tabular}{P{1.1cm}P{7.7cm}P{2.8cm}P{3.3cm}}
		\toprule
		Clave & Objetivo & Característica & Responsable \\
		\midrule
		OC-1 & \textbf{Rendimiento.} Validar a un estudiante con el prototipo en menos tiempo que con la revisión manual actual. & Eficiencia de desempeño & Ian Alexis Flores Martínez \\
		OC-2 & \textbf{Lectura confiable.} Reconocer el código al primer intento en las condiciones reales del acceso y mantener el torniquete funcionando toda la jornada. & Fiabilidad & Daniel Arellano Palacios (torniquete) e Ian Alexis Flores Martínez (lectura) \\
		OC-3 & \textbf{Seguridad y protección de datos.} Que solo un código vigente y auténtico abra el torniquete y que los datos de los estudiantes viajen y se guarden protegidos. & Seguridad & Miguel Ángel Jurado Delgadillo \\
		OC-4 & \textbf{Registro íntegro y reglas de negocio.} Registrar cada intento de forma completa, aplicar las reglas de acceso y lograr que los reportes del panel coincidan con lo ocurrido en el acceso. & Adecuación funcional & Juan Fernando Reyes Arguello \\
		OC-5 & \textbf{Facilidad de uso.} Que los estudiantes pasen sin detener el paso y que el personal consulte el panel sin ayuda. & Capacidad de interacción (usabilidad) & Bryan Uriel Villegas Raymundo \\
		\bottomrule
	\end{tabular}
	\tablenote{Fuente: Plan de Calidad (Entregable 2). Bryan Uriel Villegas Raymundo coordina la ejecución de todo el plan y Juan Esteban Anzures Campos, líder del proyecto, aprueba los resultados.}
\end{table}

\subsection{Criterios de Aceptación}
Un objetivo se da por cumplido cuando se cumplen todos sus criterios. Los que
dicen «línea base» se comparan con lo medido en el diagnóstico de octubre
(Tablas~\ref{tab:ac1} a~\ref{tab:ac2}).

\begin{table}[tbp]
	\caption{Criterios de Aceptación de OC-1 y OC-2}
	\label{tab:ac1}
	\small\setlength{\tabcolsep}{3pt}
	\begin{tabular}{P{1.2cm}P{2cm}P{7.6cm}P{4cm}}
		\toprule
		Clave & Tema & Criterio de aceptación & Evidencia \\
		\midrule
		AC-1.1 & Rendimiento & Desde que el lector detecta el código hasta que libera el torniquete pasan 3 segundos o menos en el 90\,\% de las lecturas válidas. & Registro de medición (marcas de tiempo del lector y del servidor) \\
		AC-1.2 & Rendimiento & El tiempo promedio de validación por estudiante durante el piloto es menor que el de la línea base. & Hojas de observación del flujo \\
		AC-1.3 & Rendimiento & La página del estudiante y las pantallas principales del panel cargan en menos de 3 segundos. & Registro de medición \\
		AC-2.1 & Lectura & Al menos 95\,\% de lecturas exitosas al primer intento en la matriz de condiciones de luz, brillo, distancia y modelo de celular. & Reporte de pruebas de lectura \\
		AC-2.2 & Lectura & Durante el piloto, los rechazos por error de lectura no superan el 5\,\% de los intentos. & Registros de la base de datos \\
		AC-2.3 & Lectura & Si se pierde la red, el lector lo indica en 5 segundos o menos, no abre el torniquete y conserva el intento para enviarlo al reconectarse. & Caso de prueba \\
		AC-2.4 & Lectura & El torniquete completa en banco el doble de los ciclos diarios estimados en el diagnóstico sin fallas mecánicas. & Registro de ciclos \\
		\bottomrule
	\end{tabular}
	\tablenote{Fuente: Plan de Calidad (Entregable 2).}
\end{table}

\begin{table}[tbp]
	\caption{Criterios de Aceptación de OC-3}
	\label{tab:ac1b}
	\small\setlength{\tabcolsep}{3pt}
	\begin{tabular}{P{1.2cm}P{2cm}P{7.6cm}P{4cm}}
		\toprule
		Clave & Tema & Criterio de aceptación & Evidencia \\
		\midrule
		AC-3.1 & Seguridad & Se rechaza el 100\,\% de los códigos vencidos o de las capturas reenviadas. & Caso de prueba y captura del mensaje \\
		AC-3.2 & Seguridad & Se rechaza el 100\,\% de los códigos alterados (firma HMAC inválida). & Caso de prueba y captura del mensaje \\
		AC-3.3 & Seguridad & Toda la comunicación entre el lector, el servidor y el panel usa HTTPS (conexión cifrada). & Revisión técnica (captura de tráfico) \\
		AC-3.4 & Seguridad & Las contraseñas se guardan cifradas con una función de hash (bcrypt o Argon2) y la clave secreta de la firma no sale del servidor. & Revisión técnica del código y de la base de datos \\
		AC-3.5 & Seguridad & El aviso de privacidad se muestra antes del primer uso y la base de datos guarda solo los datos necesarios. & Revisión de la propuesta con Control Escolar \\
		\bottomrule
	\end{tabular}
	\tablenote{Fuente: Plan de Calidad (Entregable 2).}
\end{table}

\begin{table}[tbp]
	\caption{Criterios de Aceptación de OC-4 y OC-5}
	\label{tab:ac2}
	\small\setlength{\tabcolsep}{3pt}
	\begin{tabular}{P{1.2cm}P{2cm}P{7.6cm}P{4cm}}
		\toprule
		Clave & Tema & Criterio de aceptación & Evidencia \\
		\midrule
		AC-4.1 & Reglas de negocio & El 100\,\% de los intentos, aceptados y rechazados, queda registrado con fecha, hora, estudiante, lector y resultado. & Caso de prueba \\
		AC-4.2 & Reglas de negocio & El torniquete solo se abre con una validación positiva; un estudiante dado de baja es rechazado y un mismo código no se acepta dos veces. & Caso de prueba y captura del mensaje \\
		AC-4.3 & Reglas de negocio & Durante el piloto, la diferencia entre el conteo manual del observador y los ingresos aceptados en la base de datos es de 2\,\% o menos. & Hoja de observación frente al reporte del panel \\
		AC-4.4 & Reglas de negocio & Los reportes del panel (ingresos por franja horaria, intentos rechazados y estudiantes sin ingresos) coinciden con consultas directas a la base de datos. & Caso de prueba \\
		AC-5.1 & Usabilidad & Al menos 85\,\% de los estudiantes del piloto pasa al primer intento sin ayuda del personal. & Hoja de observación y registros \\
		AC-5.2 & Usabilidad & Después de una explicación de 10 minutos o menos, el personal de vigilancia consulta en el panel los ingresos del día sin ayuda. & Prueba de usabilidad con tareas \\
		AC-5.3 & Usabilidad & Los estudiantes del piloto califican la facilidad de uso con un promedio de 4 o más en una escala del 1 al 5. & Encuesta breve de salida del piloto \\
		\bottomrule
	\end{tabular}
	\tablenote{Fuente: Plan de Calidad (Entregable 2).}
\end{table}

\subsection{Pruebas}
Las pruebas avanzan de lo pequeño a lo completo, siguiendo las etapas del
cronograma (Tabla~\ref{tab:niveles}). Después de cada corrección se vuelven a
ejecutar los casos afectados.

\begin{table}[tbp]
	\caption{Niveles de Prueba}
	\label{tab:niveles}
	\small\setlength{\tabcolsep}{3pt}
	\begin{tabular}{P{2.4cm}P{5.4cm}P{3.4cm}P{3.7cm}}
		\toprule
		Nivel & Qué se prueba & Fechas & Responsable \\
		\midrule
		Unitarias & Generación y validación de códigos, servicios del servidor y consultas del panel. & 26 de octubre al 27 de noviembre, en cada iteración & Miguel Ángel Jurado Delgadillo y Juan Fernando Reyes Arguello \\
		De banco & Lectura y rapidez del lector; ciclos y alimentación del torniquete. & 26 de octubre al 20 de noviembre & Ian Alexis Flores Martínez y Daniel Arellano Palacios \\
		Integración y sistema & Flujo completo en el salón: código, lector, servidor, torniquete y panel. & 23 de noviembre al 4 de diciembre & Bryan Uriel Villegas Raymundo, con todo el equipo \\
		Aceptación (piloto) & Funcionamiento en el acceso real con estudiantes del turno vespertino. & 7 al 18 de diciembre & Bryan Uriel Villegas Raymundo y Juan Esteban Anzures Campos \\
		\bottomrule
	\end{tabular}
	\tablenote{Fuente: Plan de Calidad (Entregable 2).}
\end{table}

Los 16 casos de prueba se resumen en las Tablas~\ref{tab:cp1} a~\ref{tab:cp2}.

\begin{table}[tbp]
	\caption{Casos de Prueba CP-01 a CP-04}
	\label{tab:cp1}
	\small\setlength{\tabcolsep}{3pt}
	\begin{tabular}{P{1.2cm}P{2.9cm}P{4.6cm}P{4.3cm}P{1.6cm}}
		\toprule
		Clave & Caso & Procedimiento & Resultado esperado & Criterio \\
		\midrule
		CP-01 & Ingreso válido & Un estudiante activo muestra un código vigente. & El torniquete se abre, se enciende la luz verde y el ingreso queda registrado como aceptado. & AC-4.1, AC-4.2 \\
		CP-02 & Tiempo de validación & Se hacen 50 lecturas válidas y se calculan los tiempos con las marcas del lector y del servidor. & El 90\,\% de las lecturas tarda 3 segundos o menos. & AC-1.1 \\
		CP-03 & Condiciones de lectura & 10 intentos por combinación de 3 niveles de luz, 2 de brillo, 3 distancias (15, 25 y 35 cm) y al menos 4 modelos de celular. & Al menos 95\,\% de lecturas al primer intento. & AC-2.1 \\
		CP-04 & Captura reenviada & Se toma una captura del código y se presenta cuando ya venció. & Rechazo con el mensaje «Código vencido». & AC-3.1 \\
		\bottomrule
	\end{tabular}
	\tablenote{Fuente: Plan de Calidad (Entregable 2).}
\end{table}

\begin{table}[tbp]
	\caption{Casos de Prueba CP-05 a CP-08}
	\label{tab:cp1b}
	\small\setlength{\tabcolsep}{3pt}
	\begin{tabular}{P{1.2cm}P{2.9cm}P{4.6cm}P{4.3cm}P{1.6cm}}
		\toprule
		Clave & Caso & Procedimiento & Resultado esperado & Criterio \\
		\midrule
		CP-05 & Código alterado & Se modifica un carácter del contenido del código y se genera un QR nuevo. & Rechazo con el mensaje «Código inválido». & AC-3.2 \\
		CP-06 & Reutilización & Se presenta dos veces el mismo código dentro de su vigencia. & El segundo intento se rechaza. & AC-4.2 \\
		CP-07 & Estudiante dado de baja & Un estudiante dado de baja muestra un código generado antes de la baja. & Rechazo; el torniquete no se abre. & AC-4.2 \\
		CP-08 & Falla de red & Se desconecta el Wi-Fi durante una lectura. & El lector indica la falla en 5 segundos o menos, no abre y envía el intento al reconectarse. & AC-2.3 \\
		\bottomrule
	\end{tabular}
	\tablenote{Fuente: Plan de Calidad (Entregable 2).}
\end{table}

\begin{table}[tbp]
	\caption{Casos de Prueba CP-09 a CP-16}
	\label{tab:cp2}
	\small\setlength{\tabcolsep}{3pt}
	\begin{tabular}{P{1.2cm}P{2.9cm}P{4.6cm}P{4.3cm}P{1.6cm}}
		\toprule
		Clave & Caso & Procedimiento & Resultado esperado & Criterio \\
		\midrule
		CP-09 & Comunicación cifrada & Se captura el tráfico entre el lector, el servidor y el panel. & Solo hay tráfico HTTPS; no hay datos legibles. & AC-3.3 \\
		CP-10 & Almacenamiento de contraseñas y clave & Se revisan la tabla de usuarios y el código del lector. & Solo hay hashes; la clave de la firma no aparece fuera del servidor. & AC-3.4 \\
		CP-11 & Registro completo & Se presentan 100 intentos: 80 válidos y 20 inválidos. & Hay 100 registros con todos sus campos y el resultado correcto. & AC-4.1 \\
		CP-12 & Reportes del panel & Se cargan datos de prueba conocidos y se comparan los reportes con consultas a la base. & Los valores coinciden. & AC-4.4 \\
		CP-13 & Carga de pantallas & Se carga 10 veces cada pantalla principal y se mide el tiempo con las herramientas del navegador. & Todas cargan en menos de 3 segundos. & AC-1.3 \\
		CP-14 & Resistencia del torniquete & Se acciona el torniquete en banco de forma continua. & Completa el doble de los ciclos diarios estimados sin fallas. & AC-2.4 \\
		CP-15 & Usabilidad del panel & Una persona de vigilancia hace tres tareas: ver los ingresos del día, filtrar por horario y ver los rechazos. & Las completa sin ayuda. & AC-5.2 \\
		CP-16 & Piloto en el acceso & Durante cinco días hábiles se llena la hoja de observación del flujo y se aplica la encuesta de salida. & Se cumplen los criterios del piloto. & AC-1.2, AC-2.2, AC-4.3, AC-5.1, AC-5.3 \\
		\bottomrule
	\end{tabular}
	\tablenote{Fuente: Plan de Calidad (Entregable 2).}
\end{table}

\subsection{Criterios de Entrada y Salida de las Pruebas}
\begin{itemize}
	\item \textbf{Para iniciar las pruebas en el salón:} el lector, el servidor
	      y el panel están terminados y sus pruebas unitarias y de banco están
	      aprobadas.
	\item \textbf{Para instalar el prototipo en el acceso:} al menos 90\,\% de
	      los casos ejecutados aprueba, no hay defectos críticos abiertos, los casos
	      de seguridad (CP-04, CP-05, CP-06 y CP-09) aprueban al 100\,\% y el permiso
	      de la dirección está vigente.
	\item \textbf{Para cerrar el piloto:} se completan los cinco días hábiles de
	      observación y la encuesta de salida.
\end{itemize}

\subsection{Manejo de Defectos}
Cada defecto se registra como un aviso (\textit{issue}) en el repositorio, con
el caso que lo detectó, su severidad, el responsable y su estado
(Tabla~\ref{tab:severidad}).

\begin{table}[tbp]
	\caption{Severidad de los Defectos}
	\label{tab:severidad}
	\small\setlength{\tabcolsep}{3pt}
	\begin{tabular}{P{2cm}P{7.4cm}P{5.4cm}}
		\toprule
		Severidad & Ejemplos & Atención \\
		\midrule
		Crítica & El torniquete se abre sin validación, se aceptan códigos vencidos, se exponen datos o se pierden registros. & Se detiene el avance y se corrige en 48 horas o menos. \\
		Mayor & Fallas de lectura frecuentes o reportes incorrectos en el panel. & Se corrige dentro de la misma iteración. \\
		Menor & Mensajes poco claros o detalles visuales. & Se agrega a la lista de pendientes. \\
		\bottomrule
	\end{tabular}
	\tablenote{Fuente: Plan de Calidad (Entregable 2).}
\end{table}

\subsection{Responsables de las Actividades de Calidad}
Para que nadie revise solo su propio trabajo, las revisiones técnicas son
cruzadas (Tabla~\ref{tab:respcalidad}).

\begin{table}[tbp]
	\caption{Responsables de las Actividades de Calidad}
	\label{tab:respcalidad}
	\small\setlength{\tabcolsep}{3pt}
	\begin{tabular}{P{5.4cm}P{3.5cm}P{3cm}P{3cm}}
		\toprule
		Actividad & Realiza & Aprueba & Apoya \\
		\midrule
		Coordinar el plan, diseñar los casos y elaborar el reporte de pruebas & Bryan Uriel Villegas Raymundo & Juan Esteban Anzures Campos & Todo el equipo \\
		Pruebas unitarias del servidor y de los códigos QR & Miguel Ángel Jurado Delgadillo & Bryan Uriel Villegas Raymundo & Juan Fernando Reyes Arguello \\
		Pruebas unitarias de la base de datos y del panel & Juan Fernando Reyes Arguello & Bryan Uriel Villegas Raymundo & Miguel Ángel Jurado Delgadillo \\
		Pruebas de lectura y de tiempo de validación & Ian Alexis Flores Martínez & Bryan Uriel Villegas Raymundo & Daniel Arellano Palacios \\
		Pruebas mecánicas y de alimentación del torniquete & Daniel Arellano Palacios & Bryan Uriel Villegas Raymundo & Ian Alexis Flores Martínez \\
		Revisión técnica de seguridad & Juan Fernando Reyes Arguello (revisión cruzada) & Juan Esteban Anzures Campos & Miguel Ángel Jurado Delgadillo \\
		Revisión del aviso de privacidad & Miguel Ángel Jurado Delgadillo & Juan Esteban Anzures Campos & Control Escolar \\
		Pruebas de usabilidad y encuesta de salida & Bryan Uriel Villegas Raymundo & Juan Esteban Anzures Campos & Juan Fernando Reyes Arguello \\
		Observación del acceso (línea base y piloto) & Todo el equipo, dos observadores por día & Bryan Uriel Villegas Raymundo & Personal de vigilancia \\
		Revisión de los entregables escritos & Juan Esteban Anzures Campos & Docentes asesores & Todo el equipo \\
		\bottomrule
	\end{tabular}
	\tablenote{Fuente: Plan de Calidad (Entregable 2).}
\end{table}

\subsection{Métricas}
Las métricas M2 a M4 usan los datos de la hoja de observación del flujo en el
acceso. En sus fórmulas, $\lambda$ es la tasa de llegadas (estudiantes por
minuto), $\bar{t}$ el tiempo medio de revisión en segundos, $\mu = 60/\bar{t}$
la tasa de servicio de un punto de revisión y $c$ el número de puntos de
revisión activos (Tablas~\ref{tab:met1} y~\ref{tab:met2}).

\begin{table}[tbp]
	\caption{Métricas de Calidad M1 a M6}
	\label{tab:met1}
	\small\setlength{\tabcolsep}{3pt}
	\begin{tabular}{P{0.9cm}P{2.8cm}P{4.4cm}P{3.4cm}P{3.3cm}}
		\toprule
		Clave & Métrica & Cálculo & Meta & Frecuencia \\
		\midrule
		M1 & Tiempo de validación del prototipo & Percentil 90 de (hora de apertura $-$ hora de detección del código) & 3 s o menos & Cada iteración y cada día del piloto \\
		M2 & Tiempo medio de revisión & Promedio de los tiempos cronometrados por intervalo & Menor en el piloto que en la línea base & Cada día de observación \\
		M3 & Utilización del acceso & $\rho = \lambda / (c\mu)$ en el horario pico & Menor en el piloto que en la línea base & Cada día de observación \\
		M4 & Fila máxima & Máximo de personas en fila por intervalo de 5 minutos & Menor en el piloto que en la línea base & Cada día de observación \\
		M5 & Lectura al primer intento & Lecturas exitosas al primer intento $\div$ intentos $\times$ 100 & 95\,\% o más en banco; 85\,\% o más en el piloto & Pruebas de lectura y piloto \\
		M6 & Rechazos por error de lectura & Rechazos por lectura fallida $\div$ intentos $\times$ 100 & 5\,\% o menos & Cada día del piloto \\
		\bottomrule
	\end{tabular}
	\tablenote{Fuente: Plan de Calidad (Entregable 2).}
\end{table}

\begin{table}[tbp]
	\caption{Métricas de Calidad M7 a M12}
	\label{tab:met2}
	\small\setlength{\tabcolsep}{3pt}
	\begin{tabular}{P{0.9cm}P{2.8cm}P{4.4cm}P{3.4cm}P{3.3cm}}
		\toprule
		Clave & Métrica & Cálculo & Meta & Frecuencia \\
		\midrule
		M7 & Rechazo de códigos inválidos & Códigos inválidos rechazados $\div$ códigos inválidos presentados $\times$ 100 & 100\,\% & Pruebas de seguridad \\
		M8 & Completitud del registro & Registros guardados $\div$ intentos realizados $\times$ 100 & 100\,\% & Pruebas en el salón y piloto \\
		M9 & Discrepancia de conteo & $|$conteo manual $-$ ingresos aceptados$|$ $\div$ conteo manual $\times$ 100 & 2\,\% o menos & Cada día del piloto \\
		M10 & Casos aprobados & Casos aprobados $\div$ casos ejecutados $\times$ 100 & 90\,\% o más & Al cierre de las pruebas en el salón \\
		M11 & Defectos críticos abiertos & Número de defectos críticos sin corregir & 0 antes del piloto & Revisión semanal \\
		M12 & Satisfacción de los estudiantes & Promedio de la encuesta de salida (escala del 1 al 5) & 4 o más & Al cierre del piloto \\
		\bottomrule
	\end{tabular}
	\tablenote{Fuente: Plan de Calidad (Entregable 2).}
\end{table}

\subsection{Calendario y Registros de Calidad}
Las actividades de calidad siguen el cronograma del proyecto
(Tabla~\ref{tab:calcalidad}).

\begin{table}[tbp]
	\caption{Actividades de Calidad en el Cronograma}
	\label{tab:calcalidad}
	\small\setlength{\tabcolsep}{3pt}
	\begin{tabular}{P{3.4cm}P{7.4cm}P{4cm}}
		\toprule
		Fechas & Actividad de calidad & Producto \\
		\midrule
		5 al 16 de octubre & Aprobar el plan de calidad y diseñar los casos de prueba. & Plan aprobado \\
		12 al 30 de octubre & Medir la línea base del acceso (M2, M3 y M4) durante cinco días hábiles. & Informe de diagnóstico \\
		26 de octubre al 27 de noviembre & Pruebas unitarias y de banco en cada iteración; reporte semanal de M1, M5 y M11. & Reportes semanales \\
		23 de noviembre al 4 de diciembre & Pruebas de integración y de sistema (CP-01 a CP-15) y revisión técnica de seguridad. & Reporte de pruebas y decisión de instalar \\
		7 al 18 de diciembre & Piloto en el acceso (CP-16) con medición diaria. & Hojas de observación y registros \\
		4 al 15 de enero & Comparar la línea base con el piloto y evaluar los criterios de aceptación. & Reporte final de calidad \\
		11 al 29 de enero & Documentar los resultados y las lecciones aprendidas. & Informe final \\
		\bottomrule
	\end{tabular}
	\tablenote{Fuente: Plan de Calidad (Entregable 2).}
\end{table}

Como evidencia del plan se conservarán el reporte de pruebas con el resultado
y la fecha de cada caso; los registros de medición de tiempos de validación y
de carga de pantallas; las capturas de los mensajes de rechazo y de la
revisión técnica de seguridad; las hojas de observación de la línea base y del
piloto; los avisos del repositorio con el historial de cada defecto; y las
minutas de las revisiones semanales del equipo.

\section{Seguimiento y Control}
\subsection{Cómo se Vigila el Avance}
El equipo usa cuatro mecanismos para saber si va bien:
\begin{itemize}
	\item Una junta semanal de 20 minutos, para comparar el avance con el
	      cronograma.
	\item Un tablero de tareas en el repositorio, con una columna por frente.
	\item Un reporte quincenal a la docente, con lo hecho y lo pendiente.
	\item Un control de cambios: todo lo que toque fecha, costo o alcance lo
	      autoriza el líder del proyecto y queda por escrito.
\end{itemize}
Se siguen cuatro indicadores de gestión: el porcentaje de avance de cada frente
contra lo planeado, los días de desviación respecto al cronograma por hito, la
variación del gasto acumulado con umbral de $\pm$10\,\% y los riesgos abiertos,
con cuáles ya se activaron.

\subsection{Tablero de Metas}
El tablero registra el estado de cada meta de calidad durante las pruebas y el
piloto. Todavía no hay mediciones, porque el diagnóstico empieza el 12 de
octubre. Por eso los resultados y el estado quedan pendientes
(Tabla~\ref{tab:tablero}).

\begin{table}[tbp]
	\caption{Tablero de Metas}
	\label{tab:tablero}
	\small\setlength{\tabcolsep}{3pt}
	\begin{tabular}{P{3.8cm}P{4.8cm}P{2.3cm}P{2.2cm}P{1.8cm}}
		\toprule
		Métrica & Meta y frecuencia & Resultado medido & Estado & Severidad \\
		\midrule
		M1. Tiempo de validación del prototipo & 3 segundos o menos. Cada iteración y diario en el piloto. & Por medir en las pruebas & Pendiente & Mayor \\
		M2, M3 y M4. Flujo en el acceso (tiempo medio, utilización y fila máxima) & Menores durante el piloto que en la línea base. Diario de observación. & Por medir en el diagnóstico y el piloto & Pendiente & Mayor \\
		M5 y M6. Lectura y errores (primer intento y rechazos) & Éxito de 85\,\% o más en el piloto; errores de 5\,\% o menos. Pruebas de lectura y diario en el piloto. & Por medir en las pruebas y el piloto & Pendiente & Mayor \\
		M7. Rechazo de códigos inválidos (seguridad) & 100\,\% de los códigos vencidos o alterados rechazados. Pruebas de seguridad. & Por medir en las pruebas & Pendiente & Crítica \\
		M8 y M9. Registros y conteo (adecuación funcional) & 100\,\% de registros guardados; diferencia entre el conteo manual y el sistema de 2\,\% o menos. Pruebas en el salón y diario en el piloto. & Por medir en las pruebas y el piloto & Pendiente & Crítica \\
		M11. Defectos críticos abiertos (calidad general) & 0 defectos antes de iniciar el piloto. Revisión semanal. & Por medir en las revisiones semanales & Pendiente & Crítica \\
		M12. Satisfacción estudiantil (usabilidad) & Promedio de 4 o más en una escala de 1 a 5. Al cierre del piloto, con la encuesta. & Por medir al cierre del piloto & Pendiente & Menor \\
		\bottomrule
	\end{tabular}
	\tablenote{Fuente: Tablero de Metas, Correctivos y Seguimiento (Entregable 3). En el tablero original las columnas de resultado y estado están en blanco («Cumple / No cumple»). El Entregable 3 escribe la meta de M5 como «mayor de 85\,\%»; aquí se usa «85\,\% o más» del Plan de Calidad. La métrica M10 no aparece en el tablero.}
\end{table}

\subsection{Severidad y Acciones Correctivas}
Cuando una meta no se cumple o aparece un defecto, la acción depende de su
severidad (Tabla~\ref{tab:correctivas}).

\begin{table}[tbp]
	\caption{Acciones Correctivas según la Severidad}
	\label{tab:correctivas}
	\small\setlength{\tabcolsep}{3pt}
	\begin{tabular}{P{2cm}P{5.2cm}P{4cm}P{3.6cm}}
		\toprule
		Severidad & Qué significa & Acción correctiva & Métricas del tablero \\
		\midrule
		Crítica & El torniquete se abre sin validación, se aceptan códigos vencidos, se exponen datos o se pierden registros. & Se detiene el avance y se corrige en 48 horas o menos. & M7, M8 y M9, M11 \\
		Mayor & Fallas de lectura frecuentes o reportes incorrectos en el panel. & Se corrige dentro de la misma iteración. & M1, M2 a M4, M5 y M6 \\
		Menor & Mensajes poco claros o detalles visuales. & Se agrega a la lista de pendientes. & M12 \\
		\bottomrule
	\end{tabular}
	\tablenote{Fuente: Plan de Calidad (Entregable 2) y Tablero de Metas (Entregable 3). Cada defecto se registra en el repositorio con el caso que lo detectó, su severidad, el responsable y su estado.}
\end{table}

Si un cambio afecta la fecha, el costo o el alcance, se pasa por el control de
cambios descrito arriba. Si lo que falla es una meta del tablero y no un
defecto, el equipo lo revisa en la junta semanal y decide el ajuste con el
líder del proyecto. Las fuentes no definen más pasos que estos.

\section{Metodología de Medición}
Para saber si el sistema mejora el acceso hay que medir el acceso actual y
volver a medirlo durante el piloto. El equipo usa tres instrumentos, los
mismos antes y durante el piloto, para comparar con números y no con
impresiones:
\begin{itemize}
	\item \textbf{Encuesta anónima a estudiantes} (Apéndice A del documento de
	      investigación), con preguntas cerradas. Recoge lo que solo el estudiante
	      puede decir: la espera que percibe, si olvida su credencial, si tiene
	      celular y si aceptaría mostrar un QR. La muestra es de 384 estudiantes con
	      5\,\% de error o de 196 con 7\,\%.
	\item \textbf{Entrevista a vigilancia y Control Escolar} (Apéndice B), con una
	      guía de preguntas. Aporta el procedimiento real de revisión, los casos
	      excepcionales y las condiciones para tratar datos de estudiantes.
	\item \textbf{Hoja de observación del flujo en el acceso} (Apéndice C), con
	      cronómetro y conteo, durante cinco días hábiles. Mide el tiempo de revisión
	      y la fila, que son hechos observables y no se deben preguntar.
\end{itemize}

El diseño compara un antes y un después en el mismo acceso, el mismo turno,
la misma franja horaria, los mismos observadores y el mismo instrumento. No hay
grupo de comparación, porque el acceso piloto atiende a quien llega y no es
posible operar dos accesos idénticos al mismo tiempo. Este tipo de diseño se
llama cuasiexperimental \parencite{campbellstanley1963,
hernandezsampieri2014metodologia} y ya se ha usado para evaluar un piloto
basado en código QR dentro de una institución \parencite{ortiz2026qrpilot}.
Como factores externos, por ejemplo un cambio de horario escolar, podrían
alterar los tiempos sin que el prototipo tenga que ver, los resultados se
leerán como evidencia orientativa y no concluyente. También se registra el
número de puntos de revisión activos, para no atribuir al prototipo una mejora
que venga de haber abierto más filas.

El plan de acción preliminar tiene cinco pasos: el diagnóstico con los tres
instrumentos; el diseño del hardware y de la programación del lector; el
desarrollo del servidor, la base de datos y el panel; un piloto controlado sin
retirar la revisión manual; y una evaluación que repite la observación y la
compara con la línea base.

\section{Estado Actual y Próximos Pasos}
\subsection{Lo que Ya Existe}
A la fecha de este documento (9 de octubre de 2026), las fuentes registran lo
siguiente:
\begin{itemize}
	\item El plan del proyecto se entregó el 18 de septiembre de 2026.
	\item Al 20 de septiembre estaban terminados el plan, los instrumentos de
	      medición y los avances de la investigación.
	\item El 5 de octubre de 2026, primer día del cronograma, se entregaron la
	      Ficha Inicial, el Plan de Calidad y el Tablero de Metas.
	\item El documento de investigación define el tema, el problema, la
	      pregunta, los objetivos y la metodología, con la encuesta, la entrevista y
	      la hoja de observación como apéndices.
\end{itemize}
No hay todavía mediciones del acceso, compras de materiales ni piezas del
prototipo construidas que las fuentes reporten.

\subsection{Próximos Pasos}
Según el cronograma, lo que sigue es:
\begin{enumerate}
	\item Conseguir el permiso de la dirección (16 de octubre de 2026), el
	      aviso de privacidad y la lista de materiales.
	\item Hacer el diagnóstico del 12 al 30 de octubre: encuesta, entrevistas y
	      observación del acceso durante cinco días, y compra de materiales.
	\item Pasar a las iteraciones de desarrollo a partir del 26 de octubre.
\end{enumerate}
El primer hito es el permiso de la dirección. Si no se obtiene, el plan B es
probar el prototipo en el laboratorio.

\subsection{Lo que Falta por Definir}
Estos datos no están en las fuentes y quedan pendientes:
\begin{itemize}
	\item Los resultados de todas las métricas del tablero (Tabla~\ref{tab:tablero}),
	      que dependen del diagnóstico y del piloto.
	\item La línea base del tiempo de validación, de la fila y de las incidencias
	      de identidad.
	\item El costo real, que se confirma al cotizar; el presupuesto es una
	      estimación.
	\item Los responsables de cada riesgo.
	\item Los umbrales de calidad, que son una propuesta del equipo y se ajustan
	      con la línea base del diagnóstico.
\end{itemize}


\printbibliography

\end{document}
